\chapter{Surgery for Benign Scrotal Lumps}

\section*{Introduction \& Clinical Indications}
Surgery for benign scrotal lumps encompasses a range of operative procedures aimed at the removal or repair of noncancerous masses within the scrotum. These masses can include common entities such as hydroceles (fluid collections within the tunica vaginalis), spermatoceles (epididymal cysts containing spermatozoa), epididymal cysts (simple fluid-filled cysts of the epididymis), lipomas (fatty tumors), and sebaceous cysts. The primary goal of these surgical interventions is to alleviate symptoms such as pain, discomfort, or cosmetic concerns, to confirm the benign nature of the mass, and to prevent potential complications. This surgical approach is typically considered when conservative management fails or when the mass significantly impacts the patient's quality of life, offering a definitive solution to persistent scrotal pathology.

\begin{itemize}
  \item Scrotal Mass Excision
  \item Hydrocelectomy
  \item Spermatocelectomy
  \item Epididymal Cystectomy
  \item Excision of Scrotal Cyst
  \item Tunica Vaginalis Excision or Plication
\end{itemize}

The fundamental principle underlying surgery for benign scrotal lumps is the precise identification and either removal or correction of the pathological entity while meticulously preserving the integrity and function of vital adjacent structures, particularly the testis, epididymis, and vas deferens. The rationale for intervention stems from the need to alleviate patient symptoms, confirm the benign nature of the mass, and prevent potential complications.

For hydroceles, the objective is to eliminate the fluid collection and prevent its reaccumulation. This is typically achieved by excising a portion of the tunica vaginalis sac (e.g., Jaboulay or Winkelmann technique) or by everting and plicating it (Lord's procedure) to obliterate the potential space and promote fluid absorption. In the case of spermatoceles and epididymal cysts, the surgical principle involves careful dissection and excision of the cyst from the epididymis and vas deferens. This requires a delicate technique to avoid injury to the epididymal tubules or the vas, which could compromise sperm transport and fertility. For other benign masses like lipomas or sebaceous cysts, the principle is straightforward enucleation or excision of the mass with clear margins. Across all procedures, meticulous hemostasis is crucial to prevent postoperative hematoma formation, and the ultimate technical goals include:
\begin{itemize}
  \item Complete removal or effective correction of the benign lesion.
  \item Preservation of testicular viability and function.
  \item Minimization of damage to the epididymis and vas deferens.
  \item Achievement of a satisfactory cosmetic outcome.
  \item Prevention of recurrence.
\end{itemize}

The management of scrotal swellings, particularly hydroceles, dates back to ancient civilizations, with early methods often involving simple aspiration or incision and drainage. However, these techniques were frequently associated with high recurrence rates and infection. Significant advancements began in the 18th and 19th centuries. In 1778, Percivall Pott described the anatomy of hydroceles and advocated for surgical excision. Later, in the mid-19th century, the modern surgical approaches for hydrocelectomy were refined, with techniques like the Jaboulay procedure (everting and suturing the tunica vaginalis behind the testis) gaining prominence. The Lord's procedure, involving eversion and plication of the tunica vaginalis without excision, was introduced by P. Lord in 1964, offering a less invasive alternative to complete excision, particularly for smaller hydroceles.

Spermatocele and epididymal cyst excisions evolved with improved understanding of scrotal anatomy and the development of finer surgical instruments. Early attempts were often complicated by damage to the epididymis, leading to fertility issues. The advent of microsurgical techniques in the 20th century further enhanced the precision of these delicate excisions, allowing for better preservation of the epididymis and vas deferens, thereby reducing the risk of postoperative complications and preserving fertility potential. Today, these procedures are highly refined, emphasizing minimally invasive approaches and meticulous tissue handling.

Surgery for benign scrotal lumps is indicated for a variety of reasons, primarily when the mass causes symptoms or diagnostic uncertainty.

\begin{itemize}
  \item When a scrotal lump causes persistent or increasing pain or discomfort that interferes with daily life.
  \item If the lump is significantly enlarging, causing concern, functional impairment, or difficulty with clothing.
  \item For cosmetic reasons, when the size or appearance of the lump is bothersome to the patient.
  \item When the mass interferes with daily activities, such as walking, sitting, or sexual activity.
  \item If there is diagnostic uncertainty, and malignancy cannot be definitively excluded by non-invasive means, necessitating surgical exploration.
  \item When a hydrocele is large, tense, and symptomatic, causing pressure, discomfort, or skin changes.
  \item In cases where a spermatocele or epididymal cyst is causing significant bulk, pressure symptoms, or perceived heaviness.
  \item To prevent potential complications such as rupture, infection, or testicular atrophy due to chronic compression (though rare).
\end{itemize}

Surgery for benign scrotal lumps is generally considered a primary definitive treatment for symptomatic lesions. However, the decision to proceed with surgery often follows an initial period of observation for smaller, asymptomatic masses, especially for hydroceles and spermatoceles. Scrotal ultrasound is almost always obtained as a primary diagnostic tool to confirm the benign nature of the lump, characterize its size and location, and crucially, to exclude the possibility of a testicular tumor, which would necessitate a different surgical approach (typically an inguinal orchiectomy). If a mass is confirmed benign and asymptomatic, conservative management with watchful waiting is often the initial strategy. Surgery becomes the primary intervention when symptoms develop or worsen, when the mass grows significantly, or when there is any lingering diagnostic ambiguity despite imaging. Therefore, while it is a definitive treatment, it is often initiated after non-invasive diagnostic confirmation and sometimes after a trial of conservative observation, making it a primary intervention in a carefully selected patient population.

\section*{Relevant Surgical Anatomy}
The scrotum is a fibromuscular cutaneous sac containing the testes, epididymides, and the initial portions of the spermatic cords. Its wall consists of several layers: skin, dartos muscle (responsible for scrotal rugae and temperature regulation), external spermatic fascia, cremasteric fascia (containing the cremaster muscle, which retracts the testis), and internal spermatic fascia. Deep to these layers lies the tunica vaginalis, a serous sac derived from the peritoneum, which typically has a parietal layer lining the internal spermatic fascia and a visceral layer intimately covering the testis and epididymis. A potential space, the tunica vaginalis cavity, exists between these layers and is where hydroceles form.

The testis, an ovoid organ, is suspended within the scrotum by the spermatic cord. It is responsible for spermatogenesis and hormone production. The epididymis, a C-shaped structure, lies posterolateral to the testis and consists of a head, body, and tail, serving as a site for sperm maturation and storage. The vas deferens, a muscular tube, emerges from the tail of the epididymis and ascends within the spermatic cord, transporting sperm to the ejaculatory duct. The spermatic cord also contains the testicular artery (a direct branch of the aorta), the pampiniform plexus of veins (which coalesce to form the testicular vein), lymphatic vessels draining to the para-aortic lymph nodes, and nerves including the genital branch of the genitofemoral nerve and sympathetic fibers. Surgical landmarks include the median scrotal raphe and the general contour of the testis and epididymis, which guide incision placement and dissection. Meticulous understanding of these structures is paramount to preserve testicular viability and fertility during surgery.

\section*{Preoperative Workup \& Preparation}
Prior to surgery, a thorough medical evaluation is conducted to ensure the patient is fit for anesthesia and the procedure. This comprehensive assessment aims to identify and mitigate any potential risks.

Key components of the preoperative workup and preparation include:
\begin{itemize}
  \item A detailed medical history and physical examination, focusing on cardiovascular, pulmonary, and renal systems, as well as any bleeding disorders.
  \item Blood tests, including a complete blood count, coagulation profile (PT/INR, PTT), basic metabolic panel, and sometimes liver function tests.
  \item Scrotal ultrasound to confirm the benign nature of the lump, characterize its size and location, and definitively rule out malignancy, which is crucial for surgical planning.
  \item Anesthesia consultation to assess cardiac and pulmonary fitness, discuss anesthetic options, and address any patient concerns, especially for general anesthesia.
  \item Discussion of all current medications; patients are typically advised to stop blood-thinning medications (e.g., aspirin, clopidogrel, warfarin, novel oral anticoagulants) several days before surgery, as directed by the surgeon and anesthesiologist.
  \item Patients are instructed to fast for a specified period (usually 6-8 hours for solids, 2 hours for clear liquids) before surgery to prevent aspiration during anesthesia.
  \item Prophylactic antibiotics, typically a single dose of a broad-spectrum agent (e.g., cefazolin), may be administered intravenously before the incision to reduce the risk of surgical site infection.
  \item Informed consent is obtained, ensuring the patient fully understands the procedure, its potential benefits, risks, alternatives, and the expected recovery course.
  \item Arrangements for postoperative care, including transportation home and assistance during the initial recovery period, are discussed.
\end{itemize}

\textbf{Absolute Contraindications:}
\begin{itemize}
  \item An undiagnosed scrotal mass that has not been definitively confirmed as benign, especially if a testicular tumor cannot be excluded by imaging. In such cases, an inguinal approach for exploration and possible orchiectomy is typically indicated to prevent tumor seeding into the scrotum.
  \item Severe, uncontrolled bleeding disorder or coagulopathy that cannot be corrected preoperatively, posing an unacceptable risk of hemorrhage.
  \item Patient refusal after thorough discussion of risks, benefits, and alternatives, respecting patient autonomy.
\end{itemize}

\textbf{Relative Contraindications and Precautions:}
\begin{itemize}
  \item Active scrotal infection (e.g., epididymitis, orchitis, cellulitis), which should be treated and resolved with antibiotics before elective surgery to minimize infection risk and improve surgical field clarity.
  \item Severe systemic illness or unstable medical conditions (e.g., uncontrolled diabetes, severe cardiac disease, recent myocardial infarction, severe chronic obstructive pulmonary disease) that significantly increase anesthetic or surgical risk. These conditions should be optimized prior to surgery.
  \item Mild, asymptomatic scrotal lumps where the risks of surgery (even if low) may outweigh the benefits of observation, especially in elderly or frail patients.
  \item Significant obesity, which can complicate surgical access, increase operative time, and elevate wound complication rates.
  \item History of previous scrotal surgery, which may lead to adhesions, distorted anatomy, and increased difficulty during dissection, requiring careful planning.
  \item Uncontrolled hypertension or other cardiovascular instability, which should be managed prior to elective surgery.
\end{itemize}

\section*{Surgical Technique \& Procedure}
The procedure is typically performed under general anesthesia, though regional (spinal) or local anesthesia with sedation may be used for select cases. The patient is positioned supine, often with legs slightly abducted in a frog-leg position to allow optimal access to the scrotum.

\begin{enumerate}
  \item  \textbf{Anesthesia and Positioning:} After induction of general anesthesia, the patient is positioned supine. The scrotum, perineum, and lower abdomen are prepped with an antiseptic solution and draped in a sterile fashion.
  2.  \textbf{Incision and Exposure:} A transverse scrotal incision is most commonly made, usually along a natural skin crease or the scrotal raphe, to minimize scarring. The incision is carried through the skin, dartos muscle, and superficial fascia to expose the underlying tunica vaginalis or the mass itself.
  3.  \textbf{Delivery of Testis:} The testis and epididymis are then carefully delivered through the incision, allowing for full visualization and manipulation of the scrotal contents.
  4.  \textbf{Hydrocelectomy (if applicable):} The fluid within the tunica vaginalis is aspirated. The tunica vaginalis sac is then either partially excised (Winkelmann or Jaboulay technique, where the remaining sac is everted and sutured behind the epididymis) or everted and plicated (Lord's procedure, where the sac is folded upon itself and sutured to reduce its volume and prevent fluid reaccumulation).
  5.  \textbf{Spermatocelectomy/Epididymal Cystectomy (if applicable):} The spermatocele or cyst is carefully dissected from the head or body of the epididymis, taking extreme care to preserve the epididymal tubules and the vas deferens. The stalk of the cyst, if present, is ligated, and the cyst is excised completely.
  6.  \textbf{Other Lumps (e.g., Lipoma, Sebaceous Cyst):} These are typically enucleated or excised directly from the scrotal wall or surrounding tissues, ensuring complete removal with clear margins.
  7.  \textbf{Hemostasis and Testicular Repositioning:} After the mass is removed or the hydrocele repaired, meticulous hemostasis is achieved using electrocautery or ligatures to prevent postoperative hematoma. The testis is then gently returned to the scrotum.
  8.  \textbf{Closure:} The dartos muscle layers are closed with absorbable sutures to support the testis and obliterate dead space. The skin is then closed with fine absorbable sutures, and a scrotal support dressing, often with a light compression bandage, is applied. Drains are rarely used but may be considered in cases of very large hydroceles or significant dissection where hematoma risk is high.
\end{enumerate}

\section*{Complications \& Risks}
While generally safe, surgery for benign scrotal lumps carries potential risks, which can be broadly categorized into general surgical complications and procedure-specific risks.

\begin{itemize}
  \item \textbf{General Surgical Risks:}
    \begin{itemize}
      \item \textbf{Bleeding and Hematoma:} Accumulation of blood in the scrotum, potentially requiring drainage or re-exploration.
      \item \textbf{Infection:} Surgical site infection (cellulitis, abscess), epididymitis, or orchitis, requiring antibiotic treatment or drainage.
      \item \textbf{Anesthesia Complications:} Nausea, vomiting, allergic reactions, respiratory or cardiac events, nerve injury from positioning, or adverse drug reactions.
      \item \textbf{Pain:} Postoperative pain, which is usually manageable with oral analgesics but can occasionally be severe or prolonged.
    \end{itemize}
  \item \textbf{Procedure-Specific Risks:}
    \begin{itemize}
      \item \textbf{Recurrence:} The lump (hydrocele, spermatocele, or cyst) may reaccumulate or reappear, requiring further intervention. Recurrence rates vary but can be 5-15\% for hydroceles and spermatoceles.
      \item \textbf{Injury to the Epididymis or Vas Deferens:} This is a critical risk, especially during spermatocelectomy or epididymal cystectomy, which can impair sperm transport and potentially affect fertility, particularly if bilateral or if the patient has pre-existing fertility issues.
      \item \textbf{Chronic Scrotal Pain:} Persistent pain in the scrotum, sometimes neuropathic in nature, which can be debilitating and difficult to treat.
      \item \textbf{Testicular Atrophy:} Rare, but possible if the blood supply to the testis is compromised during dissection or due to severe postoperative hematoma.
      \item \textbf{Seroma Formation:} Collection of clear, serous fluid in the surgical site, which may require aspiration or prolonged drainage.
      \item \textbf{Wound Dehiscence:} Separation of the surgical wound edges, potentially requiring re-suturing or wound care.
      \item \textbf{Cosmetic Dissatisfaction:} Scarring, asymmetry of the scrotum, or persistent swelling that may not fully resolve.
      \item \textbf{Testicular Torsion:} Extremely rare, but a theoretical risk following any scrotal manipulation, requiring immediate re-exploration.
      \item \textbf{Testicular Ischemia:} Compromise of blood flow to the testis, potentially leading to necrosis, though very uncommon with careful technique.
    \end{itemize}
\end{itemize}

\section*{Postoperative Care \& Recovery}
Immediately after the procedure, the patient is transferred to the Post-Anesthesia Care Unit (PACU) for close monitoring of vital signs, pain levels, and recovery from anesthesia. Once stable and awake, patients are typically discharged home the same day, as these are usually outpatient procedures. Oral pain medication is prescribed, and patients are advised to apply ice packs to the scrotum intermittently for the first 24-48 hours to reduce swelling and discomfort. A scrotal support or snug underwear is recommended to provide compression and support, minimizing swelling and hematoma formation. Patients should avoid strenuous activity and heavy lifting during this initial period.

During the first 1-2 weeks, patients should continue to limit strenuous activities, heavy lifting (typically anything over 10-15 lbs), and sexual activity to promote optimal healing and prevent complications. Swelling and bruising of the scrotum are common and gradually subside over this period. The incision site should be kept clean and dry, and patients are usually advised to avoid baths, opting for showers instead, to prevent wound maceration. Most individuals can return to light work or daily activities within a few days to a week, with full resumption of strenuous activities, including sports and sexual intercourse, typically allowed after 4-6 weeks, depending on the extent of surgery and individual healing. Long-term, patients are monitored for symptom resolution and absence of recurrence, with most experiencing a complete return to normal function.

During surgery, the surgeon meticulously documents the appearance, size, and location of the scrotal lump, its relationship to the testis, epididymis, and vas deferens, and the presence of any associated fluid or inflammation. For hydroceles, the findings typically include a tense, translucent, or opaque tunica vaginalis sac containing serous, straw-colored fluid. Spermatoceles are usually smooth, cystic masses arising from the epididymal head, often containing milky fluid with spermatozoa. Epididymal cysts are similar but contain clear fluid without spermatozoa. Other benign masses like lipomas appear as soft, yellowish, lobulated fatty tissue, while sebaceous cysts are typically firm, subcutaneous, and contain cheesy, malodorous material.

The resected tissue is sent for histopathological examination. The pathology report provides definitive confirmation of the benign nature of the lump. For hydroceles, the report typically describes fibrous tissue with a mesothelial lining and serous fluid. Spermatoceles are characterized by multiloculated cysts lined by cuboidal epithelium, containing spermatozoa and proteinaceous debris. Epididymal cysts show similar histology but lack spermatozoa. Lipomas are composed of mature adipose tissue, and sebaceous cysts are lined by stratified squamous epithelium and contain keratinous material. This pathological confirmation is crucial for patient reassurance and to definitively rule out any underlying malignancy.

A normal and successful postoperative course is characterized by several key findings and a predictable timeline of recovery. Patients can expect mild to moderate pain for the first few days, which is typically well-controlled with prescribed oral analgesics and gradually diminishes. Scrotal swelling and bruising are common and expected, gradually resolving over 2 to 4 weeks. The surgical incision should heal cleanly, with minimal scarring, and without signs of infection such as excessive redness, warmth, pus, or dehiscence. Upon follow-up examination, the scrotum should be soft, non-tender, and symmetric, with the testis feeling normal in size and consistency. Crucially, there should be no residual or recurrent palpable mass or significant persistent swelling. The patient should experience complete resolution of their preoperative symptoms, such as pain, pressure, or cosmetic concerns, allowing for a full return to normal activities and quality of life.

Postoperative follow-up is essential to monitor healing, ensure a successful outcome, and address any concerns.

\begin{itemize}
  \item A first postoperative visit is typically scheduled within 1-2 weeks to assess wound healing, remove any non-absorbable sutures or staples if present, and discuss the pathology results.
  \item Patients are advised on gradual resumption of physical activities, with full activity usually permitted after 4-6 weeks, including strenuous exercise and sexual activity.
  \item Instructions are provided regarding warning signs that warrant immediate medical attention, such as increasing pain, fever (temperature over 101.5$^{\circ}$F or 38.6$^{\circ}$C), significant redness or swelling, pus or foul-smelling discharge from the incision, or difficulty urinating.
  \item Further imaging, such as a repeat scrotal ultrasound, is generally not routinely performed unless there is suspicion of recurrence, a new mass, or persistent unexplained symptoms.
  \item Physical therapy or rehabilitation is typically not required for these procedures, as recovery primarily involves rest and gradual activity resumption.
\end{itemize}
Patients are encouraged to perform regular scrotal self-examinations and report any new concerns or changes to their healthcare provider.

\section*{Outcomes \& Prognosis}
Benign scrotal lumps are a common presentation in urological practice, with varying incidences depending on the specific type of lesion. Hydroceles are among the most frequent, affecting approximately 1\% of adult males, with a higher incidence in infants (congenital hydroceles often resolve spontaneously). Acquired hydroceles in adults can result from trauma, infection, inflammation, or idiopathic causes, and their incidence increases with age. Spermatoceles and epididymal cysts are also very common, found in up to 30\% of men during routine physical examination or ultrasound, particularly in middle-aged and older men. They are often asymptomatic but can grow to a size that causes discomfort. Lipomas and sebaceous cysts are less common in the scrotum specifically but can occur anywhere on the body. The decision for surgical intervention is driven by symptoms rather than prevalence. Mortality associated with surgery for benign scrotal lumps is exceedingly rare, primarily linked to severe anesthesia complications in high-risk patients. Morbidity is generally low, with the most common issues being temporary swelling, bruising, and the potential for recurrence.

The overall outcomes and prognosis for surgery for benign scrotal lumps are excellent, with high rates of symptom resolution and patient satisfaction. For hydroceles, success rates for symptom relief and prevention of recurrence are typically above 90\% with modern surgical techniques, such as the Jaboulay or Lord's procedure. Spermatoceles and epididymal cysts also have high success rates for complete excision and symptom resolution, generally exceeding 85-90\%. Recurrence rates vary by lesion type and surgical technique; hydroceles may recur in 5-10\% of cases, while spermatoceles can recur in 5-15\%, often due to incomplete excision or the development of new cysts. Functional outcomes are generally very good, with minimal long-term impact on testicular function or fertility, provided that the epididymis and vas deferens are meticulously preserved during dissection. Prognostic factors for a favorable outcome include precise surgical technique, complete removal of the lesion, and appropriate postoperative care. While complications like chronic pain or testicular atrophy are rare, they can significantly impact individual prognosis. Most patients experience a full recovery and return to normal activities within a few weeks, with a lasting resolution of their symptoms.

\section*{References}
\begin{enumerate}
  \item MedlinePlus. Scrotal Lumps. U.S. National Library of Medicine; 2024. Available at: \url{https://medlineplus.gov/scrotallumps.html}
  \item Wein AJ, Kavoussi LR, Partin AW, Peters CA, editors. Campbell-Walsh-Wein Urology. 12th ed. Philadelphia, PA: Elsevier; 2024.
  \item American Urological Association. AUA Guidelines. Available at: \url{https://www.auanet.org/guidelines-and-quality/guidelines}
  \item Netter FH. Atlas of Human Anatomy. 8th ed. Philadelphia, PA: Elsevier; 2023.
\end{enumerate}

\clearpage
